\documentclass{article}
\usepackage[T1]{fontenc}
\usepackage{lmodern}
\usepackage{graphicx}
\usepackage{amsmath,amssymb}
\usepackage[a4paper,margin=2cm]{geometry}
\usepackage[absolute,overlay]{textpos}
\setlength{\TPHorizModule}{1mm}
\setlength{\TPVertModule}{1mm}
\usepackage[indonesian]{babel}
\usepackage{hyperref}
\setlength{\emergencystretch}{2em}
% unit: Halaman 39

\begin{document}

% === Halaman 39 (python/native) ===

• Dua pihak yang berkomunikasi menyepakati parameter data sebagai 

berikut:

 1. Persamaan kurva eliptik y2  x3 + ax + b  (mod p)


- Nilai a dan b


- Bilangan prima p

 2.  Grup eliptik yang dihitung dari persamaan kurva eliptik

   3.  Titik basis (base point) B (xB, yB) , dipilih dari grup eliptik untuk operasi 

kriptografi. 

• Setiap pengguna membangkitkan pasangan kunci publik dan kunci privat

• Kunci privat = integer x, dipilih dari selang [1, p – 1]

• Kunci publik = titik Q, adalah hasil kali antara x dan titik basis B: Q = x B 

39

*) Sumber bahan: Debdeep Mukhopadhyay, Elliptic Curve Cryptography ,

 Dept of Computer Sc and Engg IIT Madras

\end{document}
